\documentclass[xcolor={table,dvipsnames},aspectratio=169,11pt]{beamer}

% --- UNIVERSAL PREAMBLE BLOCK ---
\usepackage{fontspec}
\usepackage[english, bidi=basic, provide=*]{babel}

\babelprovide[import, onchar=ids fonts]{english}

% Set default/Latin font to Sans Serif in the main (rm) and sans (sf) slots
\babelfont{rm}{Noto Sans}
\babelfont{sf}{Noto Sans}
\setsansfont{Noto Sans}

% --- SMART PACKAGES ---
\usepackage{booktabs}
\usepackage{array}
\usepackage{amssymb}
\usepackage{tikz}
\usepackage[most]{tcolorbox}

\usetikzlibrary{shapes.geometric, arrows.meta, positioning}

% --- COLOR PALETTE DEFINITIONS ---
\definecolor{nsoBlue}{HTML}{003366}
\definecolor{nsoGold}{HTML}{B48C28}
\definecolor{nsoDark}{HTML}{1A2742}
\definecolor{nsoLight}{HTML}{EDF4FC}
\definecolor{nsoMid}{HTML}{6490BE}
\definecolor{warnBg}{HTML}{FFF8E6}
\definecolor{badBg}{HTML}{FFF0F0}
\definecolor{badBorder}{HTML}{CC3333}
\definecolor{goodBg}{HTML}{F0FFF4}
\definecolor{goodBorder}{HTML}{1A7A40}

% --- BEAMER STYLE CONFIGURATION ---
\usetheme{default}
\usecolortheme{default}
\setbeamertemplate{navigation symbols}{} % Hide navigation icons

% Color mappings
\setbeamercolor{frametitle}{fg=nsoBlue}
\setbeamerfont{frametitle}{size=\large,series=\bfseries}
\setbeamercolor{normal text}{fg=nsoDark}
\setbeamercolor{structure}{fg=nsoBlue}

% Itemize template customization
\setbeamertemplate{itemize item}{\color{nsoGold}$\bullet$}
\setbeamertemplate{itemize subitem}{\color{nsoMid}$\circ$}

% Custom Frametitle with gold bottom rule
\setbeamertemplate{frametitle}{
    \vspace{0.4cm}
    \begin{minipage}{\textwidth}
        {\usebeamerfont{frametitle}\textcolor{nsoBlue}{\insertframetitle}}
        \vskip 4pt
        {\color{nsoGold}\hrule height 1.5pt}
    \end{minipage}
}

% Custom Footline (excludes title slide)
\setbeamertemplate{footline}{
    \ifnum\insertframenumber>1
    \begin{beamercolorbox}[wd=\paperwidth,dp=1ex,leftskip=0.8cm,rightskip=0.8cm]{footline}
        \color{nsoMid}\tiny NSO Nepal \hfill \insertframenumber{} / \inserttotalframenumber
    \end{beamercolorbox}
    \vspace{0.2cm}
    \fi
}

% --- CUSTOM BOXES AND BLOCKS ---
% Takeaway Box
\newtcolorbox{takeawaybox}{
    colback=nsoLight,
    colframe=nsoBlue,
    arc=4pt,
    boxrule=1pt,
    left=8pt, right=8pt, top=6pt, bottom=6pt
}

% Warning Box
\newtcolorbox{warnbox}{
    colback=warnBg,
    colframe=nsoGold,
    arc=4pt,
    boxrule=1pt,
    left=8pt, right=8pt, top=6pt, bottom=6pt
}

% Good Practice Box
\newtcolorbox{goodbox}{
    colback=goodBg,
    colframe=goodBorder,
    arc=0pt,
    boxrule=0pt,
    leftrule=3.5pt,
    left=8pt, right=8pt, top=6pt, bottom=6pt
}

% Bad Practice Box
\newtcolorbox{badbox}{
    colback=badBg,
    colframe=badBorder,
    arc=0pt,
    boxrule=0pt,
    leftrule=3.5pt,
    left=8pt, right=8pt, top=6pt, bottom=6pt
}

% Blockquote Box
\newtcolorbox{blockquote}{
    colback=nsoLight,
    colframe=nsoBlue,
    arc=0pt,
    boxrule=0pt,
    leftrule=4pt,
    left=8pt, right=8pt, top=6pt, bottom=6pt
}

% Step Badge Macro
\newcommand{\stepbadge}[1]{%
    \begin{tcolorbox}[colback=nsoGold,colframe=nsoGold,arc=8pt,auto outer arc,
        halign=center,valign=center,width=2.5cm,boxrule=0pt,top=2pt,bottom=2pt,left=2pt,right=2pt]
        \color{white}\tiny\bfseries #1
    \end{tcolorbox}%
}

% Inline Chip Macro
\newtcbox{\chip}[1][nsoBlue]{
    on line,
    arc=4pt,
    colback=#1,
    colframe=#1,
    boxrule=0pt,
    top=1.5pt,bottom=1.5pt,left=4pt,right=4pt,
    coltext=white,
    fontupper=\tiny\bfseries
}

\renewcommand{\arraystretch}{1.1}

\begin{document}

% ==========================================
% COVER SLIDE
% ==========================================
{
\setbeamercolor{background canvas}{bg=nsoBlue}
\setbeamercolor{normal text}{fg=white}
\setbeamercolor{structure}{fg=white}
\begin{frame}[plain]
    \vspace{0.8cm}
    \textcolor{nsoGold}{\rule{3cm}{4pt}}
    \vskip 0.4cm
    {\color{white}\huge\bfseries Planning and Structuring of\\Official Report Writing}
    \vskip 0.2cm
    {\color{nsoLight}\Large A Learning Guide for NSO Nepal Staff}
    \vskip 0.4cm
    \textcolor{nsoGold}{\rule{3cm}{4pt}}
    \vskip 0.4cm
    {\color{white}\bfseries National Statistics Office (NSO) Nepal} \\
    {\color{nsoLight}\small Government of Nepal}
    \vskip 0.8cm
    {\color{nsoMid}\tiny 7 Steps · Standard Depth · June 2025 · For Internal Distribution}
\end{frame}
}

% ==========================================
% AGENDA SLIDE
% ==========================================
{
\setbeamercolor{background canvas}{bg=nsoLight}
\begin{frame}{Agenda}
    \rowcolors{2}{white}{nsoLight!50}
    \centering
    \small
    \begin{tabular}{cl}
    \toprule
    \rowcolor{nsoBlue}\color{white}\bfseries Step & \color{white}\bfseries Topic \\
    \midrule
    \textbf{1} & The Landscape of Official Reports \\
    \textbf{2} & Choosing the Right Report for the Right Situation \\
    \textbf{3} & Structuring the Outline \\
    \textbf{4} & What to Include in Each Section \\
    \textbf{5} & Writing with Statistical Precision \\
    \textbf{6} & Review, Quality Assurance \& Finalization \\
    \textbf{7} & Synthesis: From Blank Page to Published Report \\
    $\cdot$ & Summary \& Next Steps \\
    \bottomrule
    \end{tabular}
\end{frame}
}

% ==========================================
% STEP 1 HEADER
% ==========================================
{
\setbeamercolor{background canvas}{bg=nsoBlue}
\setbeamercolor{normal text}{fg=white}
\setbeamercolor{structure}{fg=white}
\begin{frame}[plain]
    \vspace{1.5cm}
    \stepbadge{Step 1 of 7}
    \vskip 0.4cm
    {\color{white}\LARGE\bfseries The Landscape of Official Reports}
    \vskip 0.4cm
    {\color{nsoLight}\large Understanding the types of official and statistical reports, their purposes, and how they differ from general writing.}
    \vskip 0.8cm
    \textcolor{nsoGold}{\rule{3cm}{3pt}}
\end{frame}
}

% ==========================================
% WHAT IS AN OFFICIAL REPORT?
% ==========================================
\begin{frame}{What Is an ``Official Report''?}
    An official report is a \textbf{structured, purposeful document} that communicates factual information to a defined audience. Unlike general writing, official reports must be:
    \vskip 0.2cm
    \rowcolors{2}{nsoLight}{white}
    \footnotesize
    \begin{tabular}{lp{10cm}}
    \toprule
    \rowcolor{nsoBlue}\color{white}\bfseries Property & \color{white}\bfseries Meaning \\
    \midrule
    \textbf{Accurate} & Numbers and facts are verified \\
    \textbf{Objective} & Free from personal opinion \\
    \textbf{Neutral in tone} & Descriptive, not persuasive \\
    \textbf{Traceable} & Claims backed by data or citations \\
    \bottomrule
    \end{tabular}
    \vskip 0.2cm
    \begin{blockquote}
    \footnotesize For NSO Nepal, reports are the primary channel for sharing national statistics with government, media, researchers, and international partners.
    \end{blockquote}
\end{frame}

% ==========================================
% 8 TYPES OF REPORTS
% ==========================================
\begin{frame}{8 Types of Official and Statistical Reports}
    \rowcolors{2}{nsoLight}{white}
    \tiny
    \begin{tabular}{ll>{\raggedright}p{5.5cm}>{\raggedright\arraybackslash}p{4cm}}
    \toprule
    \rowcolor{nsoBlue}\color{white}\bfseries \# & \color{white}\bfseries Report Type & \color{white}\bfseries Core Purpose & \color{white}\bfseries Primary Audience \\
    \midrule
    1 & Statistical Bulletin & Publish key indicators (CPI, GDP) & Government, media, public \\
    2 & Press Release & Announce findings quickly \& accessibly & Media, general public \\
    3 & Statistical Abstract & Comprehensive data across sectors & Researchers, planners \\
    4 & Analytical / Thematic Report & In-depth topic analysis & Policymakers, donors \\
    5 & Census / Survey Report & Full methodology \& findings & All stakeholders \\
    6 & Administrative Report & Progress updates, minutes & NSO management \\
    7 & Technical / Methodology Report & How data was collected \& processed & Statisticians \\
    8 & Policy Brief & Action-oriented summary & Ministers, officials \\
    \bottomrule
    \end{tabular}
\end{frame}

% ==========================================
% KEY REPORT TYPES AT A GLANCE
% ==========================================
\begin{frame}{Key Report Types --- At a Glance}
    \begin{columns}[t]
        \begin{column}{0.48\textwidth}
            \small
            \textbf{Statistical Bulletin} \\
            Fixed-schedule publication (monthly, quarterly, annually). Numbers, tables, brief commentary. \\
            \textit{Example: Nepal CPI --- Monthly Bulletin}
            \vskip 0.2cm
            \textbf{Press Release} \\
            1--2 page announcement for journalists \& public. Simple, direct, jargon-free. First public-facing document alongside any major release.
            \vskip 0.2cm
            \textbf{Census / Survey Report} \\
            Most comprehensive type. Covers background, methodology, findings, annexes. \\
            \textit{Example: NPHC 2078 --- National Report}
        \end{column}
        \begin{column}{0.48\textwidth}
            \small
            \textbf{Analytical / Thematic Report} \\
            Explains \emph{what the numbers mean}. Includes context, comparisons, sometimes recommendations.
            \vskip 0.2cm
            \textbf{Policy Brief} \\
            2--4 pages. Explicitly recommends action. A Minister should be able to read it in \textbf{5 minutes}.
            \vskip 0.2cm
            \textbf{Technical / Methodology Report} \\
            For statisticians and peer institutions. Explains sampling design, processing steps, caveats.
        \end{column}
    \end{columns}
\end{frame}

% ==========================================
% WHY DOES THE DISTINCTION MATTER?
% ==========================================
\begin{frame}{Why Does the Distinction Matter?}
    Each report type has \textbf{different rules} for:
    \begin{itemize}
        \item \textbf{Length} --- from 1 page (press release) to 300+ pages (census report)
        \item \textbf{Tone} --- technical vs. plain language
        \item \textbf{Structure} --- free-form vs. fixed template
        \item \textbf{Audience} --- statisticians vs. journalists vs. ministers
    \end{itemize}
    \begin{warnbox}
    \scriptsize \textbf{Common mistakes:} Writing a press release that reads like a methodology report, or burying a policy brief's recommendation on page 10. Choosing the wrong report type wastes time and reduces impact.
    \end{warnbox}
    \begin{takeawaybox}
    \scriptsize \textbf{Key Takeaway:} The first decision in any writing task is --- \emph{What kind of report is this, and who is it for?}
    \end{takeawaybox}
\end{frame}

% ==========================================
% STEP 2 HEADER
% ==========================================
{
\setbeamercolor{background canvas}{bg=nsoBlue}
\setbeamercolor{normal text}{fg=white}
\setbeamercolor{structure}{fg=white}
\begin{frame}[plain]
    \vspace{1.5cm}
    \stepbadge{Step 2 of 7}
    \vskip 0.4cm
    {\color{white}\LARGE\bfseries Choosing the Right Report for the Right Situation}
    \vskip 0.4cm
    {\color{nsoLight}\large How to decide which report type to write based on audience, purpose, and occasion.}
    \vskip 0.8cm
    \textcolor{nsoGold}{\rule{3cm}{3pt}}
\end{frame}
}

% ==========================================
% ASK THREE QUESTIONS FIRST
% ==========================================
\begin{frame}{Ask Three Questions First}
    Before writing a single word, answer these:
    \vskip 0.4cm
    \begin{blockquote}
        \textbf{1. Who is my audience?} \\
        \small Journalists? Ministers? Statisticians? The general public?
    \end{blockquote}
    \begin{blockquote}
        \textbf{2. What do I want from them?} \\
        \small To report the news? To make a policy decision? To replicate our methodology?
    \end{blockquote}
    \begin{blockquote}
        \textbf{3. What is the urgency and occasion?} \\
        \small Scheduled release? One-time survey result? Emergency announcement?
    \end{blockquote}
    \vskip 0.2cm
    Your answers point directly to the right report type.
\end{frame}

% ==========================================
% DECISION GUIDE
% ==========================================
\begin{frame}{Decision Guide: When to Write What}
    \rowcolors{2}{nsoLight}{white}
    \footnotesize
    \begin{tabular}{l>{\raggedright\arraybackslash}p{6.5cm}}
    \toprule
    \rowcolor{nsoBlue}\color{white}\bfseries Situation & \color{white}\bfseries Best Report Type \\
    \midrule
    Monthly/quarterly data ready (CPI, trade, employment) & \textbf{Statistical Bulletin} \\
    Publishing a bulletin and want media coverage & \textbf{Press Release} \textit{(alongside bulletin)} \\
    Minister needs a 5-minute read before a meeting & \textbf{Policy Brief} \\
    Completed a national census or survey & \textbf{Census / Survey Report} \\
    Analysing trends across multiple indicators over time & \textbf{Analytical / Thematic Report} \\
    Peer institution asks how you collected your data & \textbf{Technical / Methodology Report} \\
    Briefing your director on project progress & \textbf{Administrative / Internal Report} \\
    International donor wants a country data summary & \textbf{Statistical Abstract / Compendium} \\
    \bottomrule
    \end{tabular}
\end{frame}

% ==========================================
% THE PRESS RELEASE
% ==========================================
\begin{frame}{The Press Release --- When and Why}
    \bullet A press release should be issued \textbf{every time} NSO Nepal publishes data with public significance.
    \vskip 0.2cm
    \textbf{What it must do:}
    \begin{itemize}
        \item Released \textbf{simultaneously} with the main statistical publication.
        \item Sent \textbf{directly} to journalists, news portals, social media teams.
        \item Contain \textbf{one or two headline numbers} --- not the full dataset.
        \item Written so a journalist can reproduce the story almost word for word.
    \end{itemize}
    \vskip 0.2cm
    \begin{warnbox}
    \footnotesize \textbf{Common Mistake:} Issuing a press release \emph{days after} the bulletin, when journalists have already sourced data elsewhere --- often incorrectly. \textbf{Timing is everything.}
    \end{warnbox}
\end{frame}

% ==========================================
% PRACTICAL EXAMPLE
% ==========================================
\begin{frame}{Practical Example: NLFS 2082}
    One survey $\rightarrow$ five different documents:
    \vskip 0.2cm
    \rowcolors{2}{nsoLight}{white}
    \footnotesize
    \begin{tabular}{clll}
    \toprule
    \rowcolor{nsoBlue}\color{white}\bfseries Order & \color{white}\bfseries Document & \color{white}\bfseries Type & \color{white}\bfseries Audience \\
    \midrule
    1 & Sampling \& Field Operations Report & Technical/Methodology & Statisticians (internal first) \\
    2 & NLFS 2082 Main Report & Survey Report & All stakeholders \\
    3 & Youth Unemployment --- Thematic Analysis & Analytical Report & Policymakers, donors \\
    4 & NLFS 2082 Launch Announcement & \textbf{Press Release} & Media, public \\
    5 & Labour Market Brief for Ministry & Policy Brief & Ministry of Labour \\
    \bottomrule
    \end{tabular}
    \vskip 0.2cm
    \begin{takeawaybox}
    \scriptsize \textbf{Key Takeaway:} The report type is determined by \emph{who reads it} and \emph{what they need to do with it.}
    \end{takeawaybox}
\end{frame}

% ==========================================
% STEP 3 HEADER
% ==========================================
{
\setbeamercolor{background canvas}{bg=nsoBlue}
\setbeamercolor{normal text}{fg=white}
\setbeamercolor{structure}{fg=white}
\begin{frame}[plain]
    \vspace{1.5cm}
    \stepbadge{Step 3 of 7}
    \vskip 0.4cm
    {\color{white}\LARGE\bfseries Structuring the Outline}
    \vskip 0.4cm
    {\color{nsoLight}\large How to plan a report's skeleton before writing --- and why structure matters for credibility.}
    \vskip 0.8cm
    \textcolor{nsoGold}{\rule{3cm}{3pt}}
\end{frame}
}

% ==========================================
% WHY OUTLINING COMES FIRST
% ==========================================
\begin{frame}{Why Outlining Comes First}
    Writing without an outline is like \textbf{building a house without a blueprint}.
    \vskip 0.2cm
    \textbf{A solid outline:}
    \begin{itemize}
        \item Ensures no critical section is missed.
        \item Lets multiple authors divide the work clearly.
        \item Prevents the report from growing in unplanned directions.
        \item Makes peer review and editing much faster.
    \end{itemize}
    \vskip 0.2cm
    \begin{blockquote}
    \footnotesize \textbf{Rule:} Always outline \emph{before} you write. The investment pays back tenfold in writing speed and quality.
    \end{blockquote}
\end{frame}

% ==========================================
% THE UNIVERSAL OUTLINE FRAMEWORK
% ==========================================
\begin{frame}{The Universal Outline Framework}
    \begin{columns}[t]
    \begin{column}{0.48\textwidth}
        \begin{itemize}
            \item[\textbf{1.}] \textbf{Front Matter}
            \begin{itemize}
                \item Title page
                \item Foreword / Preface
                \item Table of Contents
                \item List of Tables \& Figures
                \item Abbreviations
            \end{itemize}
            \vskip 0.1cm
            \item[\textbf{2.}] \textbf{Introductory Section}
            \begin{itemize}
                \item Background / Context
                \item Objectives
                \item Scope \& Coverage
            \end{itemize}
        \end{itemize}
    \end{column}
    \begin{column}{0.48\textwidth}
        \begin{itemize}
            \item[\textbf{3.}] \textbf{Core Content}
            \begin{itemize}
                \item Methodology
                \item Findings / Results
                \item Analysis / Interpretation
            \end{itemize}
            \vskip 0.1cm
            \item[\textbf{4.}] \textbf{Concluding Section}
            \begin{itemize}
                \item Summary / Conclusions
                \item Recommendations (where applicable)
            \end{itemize}
            \vskip 0.1cm
            \item[\textbf{5.}] \textbf{Back Matter}
            \begin{itemize}
                \item References
                \item Annexes / Appendices
                \item Glossary
            \end{itemize}
        \end{itemize}
    \end{column}
    \end{columns}
\end{frame}

% ==========================================
% OUTLINES BY REPORT TYPE
% ==========================================
\begin{frame}{Outlines by Report Type}
    \begin{columns}[t]
        \begin{column}{0.48\textwidth}
            \begin{tcolorbox}[colback=nsoLight, colframe=nsoBlue, title=Statistical Bulletin, coltitle=white, boxrule=1pt, arc=3pt, fonttitle=\scriptsize\bfseries, fontupper=\scriptsize]
                Cover $\rightarrow$ Key Highlights $\rightarrow$ \\
                Data Tables $\rightarrow$ \\
                Short Analytical Notes $\rightarrow$ \\
                Methodology Note $\rightarrow$ Contact
            \end{tcolorbox}
            \begin{tcolorbox}[colback=nsoLight, colframe=nsoBlue, title=Press Release, coltitle=white, boxrule=1pt, arc=3pt, fonttitle=\scriptsize\bfseries, fontupper=\scriptsize]
                Headline $\rightarrow$ Dateline $\rightarrow$ \\
                Opening paragraph (the news) $\rightarrow$ \\
                Key findings (2--3 paragraphs) $\rightarrow$ \\
                Context $\rightarrow$ NSO Chief Quote $\rightarrow$ \\
                Boilerplate $\rightarrow$ Contact
            \end{tcolorbox}
        \end{column}
        \begin{column}{0.48\textwidth}
            \begin{tcolorbox}[colback=nsoLight, colframe=nsoBlue, title=Analytical Report, coltitle=white, boxrule=1pt, arc=3pt, fonttitle=\scriptsize\bfseries, fontupper=\scriptsize]
                Cover $\rightarrow$ Foreword $\rightarrow$ \\
                Executive Summary $\rightarrow$ \\
                Introduction $\rightarrow$ \\
                Data Sources \& Methods $\rightarrow$ \\
                Thematic Findings $\rightarrow$ \\
                Conclusions \& Recommendations $\rightarrow$ \\
                References $\rightarrow$ Annexes
            \end{tcolorbox}
            \begin{tcolorbox}[colback=nsoLight, colframe=nsoBlue, title=Policy Brief, coltitle=white, boxrule=1pt, arc=3pt, fonttitle=\scriptsize\bfseries, fontupper=\scriptsize]
                Title $\rightarrow$ Key Message $\rightarrow$ \\
                Background $\rightarrow$ \\
                Evidence (2--3 findings) $\rightarrow$ \\
                Policy Options $\rightarrow$ \\
                Recommended Action $\rightarrow$ \\
                Further Reading
            \end{tcolorbox}
        \end{column}
    \end{columns}
\end{frame}

% ==========================================
% BUILDING YOUR OUTLINE
% ==========================================
\begin{frame}{Building Your Outline --- 5 Steps}
    \begin{enumerate}
        \item \textbf{Start with purpose} --- Write one sentence: \\ \textit{``This report will help [audience] understand/decide [X].''}
        \item \textbf{List the questions your reader will have} --- your sections should answer each one.
        \item \textbf{Group related questions} into logical chapters or sections.
        \item \textbf{Order them} --- general to specific, or context $\rightarrow$ findings $\rightarrow$ recommendations.
        \item \textbf{Assign approximate length} --- so you know where the weight falls.
    \end{enumerate}
    \vskip 0.2cm
    \begin{columns}
        \begin{column}{0.5\textwidth}
            \begin{takeawaybox}
            \scriptsize \textbf{NSO Nepal Tip:} For census/survey reports, check UNFPA, UNICEF, or UN Statistics Division templates first.
            \end{takeawaybox}
        \end{column}
        \begin{column}{0.5\textwidth}
            \begin{takeawaybox}
            \scriptsize \textbf{Key Takeaway:} An outline is not a table of contents --- it is a \emph{thinking tool.}
            \end{takeawaybox}
        \end{column}
    \end{columns}
\end{frame}

% ==========================================
% STEP 4 HEADER
% ==========================================
{
\setbeamercolor{background canvas}{bg=nsoBlue}
\setbeamercolor{normal text}{fg=white}
\setbeamercolor{structure}{fg=white}
\begin{frame}[plain]
    \vspace{1.5cm}
    \stepbadge{Step 4 of 7}
    \vskip 0.4cm
    {\color{white}\LARGE\bfseries What to Include in Each Section}
    \vskip 0.4cm
    {\color{nsoLight}\large The anatomy of a well-written official report --- section by section.}
    \vskip 0.8cm
    \textcolor{nsoGold}{\rule{3cm}{3pt}}
\end{frame}
}

% ==========================================
% FRONT MATTER SECTIONS
% ==========================================
\begin{frame}{Front Matter Sections}
    \rowcolors{2}{nsoLight}{white}
    \footnotesize
    \begin{tabular}{lp{10cm}}
    \toprule
    \rowcolor{nsoBlue}\color{white}\bfseries Section & \color{white}\bfseries What to Include \\
    \midrule
    \textbf{Title Page} & Report title, subtitle, NSO Nepal, date/period, edition, GoN emblem \\
    \textbf{Foreword / Preface} & Importance of the report, thanks to contributors, \textbf{no findings} (1--2 paragraphs) \\
    \textbf{Executive Summary} & Stand-alone summary; purpose + key findings + recommendations; \textbf{written last}; 1--2 pages \\
    \textbf{Table of Contents} & Auto-generated; must match actual section headings and page numbers \\
    \textbf{Abbreviations} & All acronyms defined in one place \\
    \bottomrule
    \end{tabular}
    \vskip 0.2cm
    \begin{takeawaybox}
    \footnotesize \textbf{Important:} The \textbf{Executive Summary} is the most-read section of any long report. A reader who reads \emph{only this} must understand the report's main message.
    \end{takeawaybox}
\end{frame}

% ==========================================
% CORE CONTENT SECTIONS
% ==========================================
\begin{frame}{Core Content Sections}
    \small
    \textbf{Introduction / Background} answers:
    \begin{itemize}
        \item Why was this report written?
        \item What time period, geography, and population does it cover?
        \item How does it connect to previous reports in the series?
    \end{itemize}
    \vskip 0.15cm
    \textbf{Methodology / Data Sources} must include:
    \begin{itemize}
        \item Data source(s): survey, census, administrative records, satellite data.
        \item Sample design: sampling frame, sample size, stratification.
        \item Data collection period and methods.
        \item Processing steps: editing, imputation, weighting.
        \item \textbf{Key limitations and caveats}.
    \end{itemize}
    \vskip 0.15cm
    \textbf{Findings / Results} --- structured by thematic chapters or indicator sections. Each finding needs: \textbf{the number $\rightarrow$ comparison $\rightarrow$ interpretation $\rightarrow$ supporting table/chart}.
\end{frame}

% ==========================================
% ANALYTICAL AND CLOSING SECTIONS
% ==========================================
\begin{frame}{Analytical and Closing Sections}
    \small
    \textbf{Analysis / Discussion} \\
    Goes beyond numbers to explain them objectively:
    \begin{itemize}
        \item What patterns emerge?
        \item Are there regional, gender, or income-level differences?
        \item How does this compare to other countries or previous surveys?
    \end{itemize}
    \vskip 0.15cm
    \textbf{Conclusions and Recommendations}
    \begin{itemize}
        \item \textbf{Conclusions} = what the data \emph{shows}.
        \item \textbf{Recommendations} = what should be \emph{done} \textit{(only when the brief calls for it)}.
    \end{itemize}
    \vskip 0.15cm
    \textbf{Annexes} \\
    House technical material that would slow down the main text: \\
    \small full data tables · survey questionnaires · sampling frame details · definitions \& classifications.
\end{frame}

% ==========================================
% THE PRESS RELEASE SECTION BY SECTION
% ==========================================
\begin{frame}[fragile]{The Press Release --- Section by Section}
    \begin{tcolorbox}[colback=white, colframe=nsoBlue, arc=4pt, title=Press Release Structure, coltitle=white, boxrule=1pt]
    \tiny
    \texttt{\textbf{[NSO NEPAL LOGO / LETTERHEAD]}} \\
    \texttt{\textbf{FOR IMMEDIATE RELEASE} \hfill Date: Baisakh 15, 2082 / 28 April 2025} \\[3pt]
    \texttt{\textbf{HEADLINE: Nepal's Inflation Rises to 6.2\% in Baisakh 2082}} \\
    \textit{\scriptsize (One sentence · factual · newsworthy)} \\[3pt]
    \texttt{KATHMANDU — [Opening paragraph: single most important fact in plain language. What, when, how much?]} \\[2pt]
    \texttt{[Body paragraph 1: Second most important finding + comparison to prior period]} \\[2pt]
    \texttt{[Body paragraph 2: Context — what is driving the change?]} \\[2pt]
    \texttt{``Quote,'' said [Name], Chief Statistician, NSO Nepal. \hfill Contact: [Name · Title · Email · Phone]} \\[2pt]
    \hfill \texttt{\textbf{\#\#\#}}
    \end{tcolorbox}
\end{frame}

% ==========================================
% THE GOLDEN RULE OF PRESS RELEASES
% ==========================================
\begin{frame}{The Golden Rule of Press Releases}
    \begin{columns}[c]
        \begin{column}{0.5\textwidth}
            \small
            \textbf{Inverted Pyramid Structure} \\
            Most important $\rightarrow$ least important, top to bottom.
            \vskip 0.2cm
            Journalists cut from the \textbf{bottom}. If trimmed, the headline + opening paragraph must still tell the complete story.
        \end{column}
        \begin{column}{0.48\textwidth}
            \centering
            \begin{tikzpicture}[scale=0.8, every node/.style={transform shape}]
                % Trapezoid 1 (Headline)
                \draw[fill=nsoBlue, draw=none] (0,3) -- (6,3) -- (5.5,2.2) -- (0.5,2.2) -- cycle;
                \node[white] at (3,2.6) {\scriptsize\bfseries HEADLINE NUMBER / FACT};
                \node[right, nsoDark] at (6,2.6) {\tiny $\leftarrow$ Most important};
                
                % Trapezoid 2 (Second key)
                \draw[fill=nsoBlue!80!nsoGold, draw=none] (0.5,2.1) -- (5.5,2.1) -- (5.0,1.4) -- (1.0,1.4) -- cycle;
                \node[white] at (3,1.75) {\scriptsize\bfseries Second key finding};
                
                % Trapezoid 3 (Context)
                \draw[fill=nsoBlue!60!nsoGold, draw=none] (1.0,1.3) -- (5.0,1.3) -- (4.5,0.7) -- (1.5,0.7) -- cycle;
                \node[white] at (3,1.0) {\scriptsize Context / background};
                
                % Trapezoid 4 (Quote)
                \draw[fill=nsoMid, draw=none] (1.5,0.6) -- (4.5,0.6) -- (4.0,0.1) -- (2.0,0.1) -- cycle;
                \node[white] at (3,0.35) {\scriptsize NSO Chief Quote};
                
                % Trapezoid 5 (Boilerplate)
                \draw[fill=nsoLight, draw=nsoBlue, line width=0.5pt] (2.0,0.0) -- (4.0,0.0) -- (3.7,-0.5) -- (2.3,-0.5) -- cycle;
                \node[nsoDark] at (3,-0.25) {\scriptsize\bfseries Boilerplate};
                \node[right, nsoDark] at (4,-0.25) {\tiny $\leftarrow$ Least important};
            \end{tikzpicture}
        \end{column}
    \end{columns}
\end{frame}

% ==========================================
% STEP 5 HEADER
% ==========================================
{
\setbeamercolor{background canvas}{bg=nsoBlue}
\setbeamercolor{normal text}{fg=white}
\setbeamercolor{structure}{fg=white}
\begin{frame}[plain]
    \vspace{1.5cm}
    \stepbadge{Step 5 of 7}
    \vskip 0.4cm
    {\color{white}\LARGE\bfseries Writing with Statistical Precision}
    \vskip 0.4cm
    {\color{nsoLight}\large How to present numbers, analysis, and language accurately and accessibly.}
    \vskip 0.8cm
    \textcolor{nsoGold}{\rule{3cm}{3pt}}
\end{frame}
}

% ==========================================
% PRINCIPLE 1 - LET THE NUMBER LEAD
% ==========================================
\begin{frame}{Principle 1 --- Let the Number Lead}
    Always give the actual number first. Never lead with vague adjectives.
    \vskip 0.2cm
    \begin{columns}[t]
    \begin{column}{0.48\textwidth}
        \begin{badbox}
        \footnotesize\textit{``There has been a \textbf{significant} upward trend in price levels.''}
        \end{badbox}
    \end{column}
    \begin{column}{0.48\textwidth}
        \begin{goodbox}
        \footnotesize\textit{``Consumer prices rose by \textbf{6.2\%} year-on-year in Baisakh 2082, up from \textbf{5.8\%} in Chaitra 2081.''}
        \end{goodbox}
    \end{column}
    \end{columns}
    \vskip 0.2cm
    \textbf{Avoid these vague words in statistical writing:}
    \vskip 0.1cm
    \rowcolors{2}{nsoLight}{white}
    \footnotesize
    \begin{tabular}{ll}
    \toprule
    \rowcolor{nsoBlue}\color{white}\bfseries Vague Word & \color{white}\bfseries Replace with... \\
    \midrule
    significant & [actual number] \\
    high / low & [number + context] \\
    considerable & [number + comparison] \\
    dramatic change & [before figure] $\rightarrow$ [after figure] \\
    \bottomrule
    \end{tabular}
\end{frame}

% ==========================================
% PRINCIPLE 2 - BE PRECISE
% ==========================================
\begin{frame}{Principle 2 --- Be Precise About What You Measure}
    Every number must be anchored in \textbf{three dimensions}:
    \vskip 0.4cm
    \rowcolors{2}{nsoLight}{white}
    \small
    \begin{tabular}{llp{6.5cm}}
    \toprule
    \rowcolor{nsoBlue}\color{white}\bfseries Dimension & \color{white}\bfseries Question to answer & \color{white}\bfseries Example \\
    \midrule
    \textbf{Indicator} & What exactly is being measured? & Unemployment \emph{rate} (not persons) \\
    \textbf{Reference period} & When? & Baisakh 2082 (Nepali) = April/May 2025 \\
    \textbf{Population} & Who is included? & Persons aged 15+, all 77 districts \\
    \bottomrule
    \end{tabular}
\end{frame}

% ==========================================
% PRINCIPLE 3 - PERCENTAGE VS PP
% ==========================================
\begin{frame}{Principle 3 --- Percentage vs. Percentage Points}
    \textbf{This is the most common error in statistical writing.}
    \vskip 0.2cm
    \begin{columns}[t]
    \begin{column}{0.48\textwidth}
        \textbf{Scenario:} Inflation moves from 5\% to 6\%. \\[6pt]
        \textbf{Change:}
        \begin{itemize}
            \item \textcolor{goodBorder}{\textbf{Correct:}} Rose by \textbf{1 percentage point}
            \item \textcolor{badBorder}{\textbf{Incorrect:}} Rose by 1\% $\leftarrow$ \textbf{WRONG}
        \end{itemize}
        \vskip 0.2cm
        A 1\% \emph{increase} in a 5\% rate is:
        $$5\% \times 1.01 = 5.05\% \quad (\text{not } 6\%)$$
    \end{column}
    \begin{column}{0.48\textwidth}
        \textbf{Why it matters:} \\[6pt]
        Getting this wrong causes journalists to misreport your figure by a \textbf{factor of 5}.
        \vskip 0.3cm
        \begin{takeawaybox}
        \footnotesize \textbf{Always say:} \\[2pt]
        \emph{``rose by X percentage points''} \\[2pt]
        not \emph{``rose by X percent''} when measuring a change in a rate.
        \end{takeawaybox}
    \end{column}
    \end{columns}
\end{frame}

% ==========================================
% PRINCIPLES 4 & 5 - TABLES & UNCERTAINTY
% ==========================================
\begin{frame}{Principles 4 \& 5 --- Tables and Uncertainty}
    \textbf{Every table and figure must have:}
    \rowcolors{2}{nsoLight}{white}
    \footnotesize
    \begin{tabular}{lp{10cm}}
    \toprule
    \rowcolor{nsoBlue}\color{white}\bfseries Element & \color{white}\bfseries Requirement \\
    \midrule
    \textbf{Number} & Table 3.1, Figure 2 \\
    \textbf{Title} & Fully self-describing (no need to read the text) \\
    \textbf{Unit} & \% · Rs · persons · thousands --- in the header \\
    \textbf{Source line} & \emph{Source: NSO Nepal, NLFS 2082} \\
    \textbf{Notes} & Footnotes for exclusions, rounding, caveats \\
    \bottomrule
    \end{tabular}
    \vskip 0.2cm
    \begin{itemize}
        \item Keep tables \textbf{in the text near where they are referenced} --- don't hide them in annexes.
        \item \textbf{Qualify uncertainty honestly} --- margins of error, small sample size warnings, and data quality flags. Transparency is a mark of credibility, not weakness.
    \end{itemize}
\end{frame}

% ==========================================
% PRECISION IN PRESS RELEASES
% ==========================================
\begin{frame}{Precision in Press Releases}
    \rowcolors{2}{nsoLight}{white}
    \small
    \begin{tabular}{lp{9cm}}
    \toprule
    \rowcolor{nsoBlue}\color{white}\bfseries Rule & \color{white}\bfseries Example \\
    \midrule
    \textbf{Round intelligently in headlines} & ``about 1.2 million persons'' not ``1,213,447'' \\
    \textbf{Give precise figures in the body} & Full number in paragraph 2 or linked table \\
    \textbf{Anchor comparisons} & ``fell from 25.2\% in 2010/11 to 17.4\% in 2022/23'' \\
    \textbf{Define technical terms} & ``labour force participation rate \textit{(the share of working-age population employed or actively seeking work)}'' \\
    \bottomrule
    \end{tabular}
    \vskip 0.2cm
    \begin{takeawaybox}
    \scriptsize \textbf{Key Takeaway:} Precision is not about complexity --- it is about \emph{exactness}. Numbers must be anchored in time, population, and units. Vague statistical writing erodes public trust.
    \end{takeawaybox}
\end{frame}

% ==========================================
% STEP 6 HEADER
% ==========================================
{
\setbeamercolor{background canvas}{bg=nsoBlue}
\setbeamercolor{normal text}{fg=white}
\setbeamercolor{structure}{fg=white}
\begin{frame}[plain]
    \vspace{1.5cm}
    \stepbadge{Step 6 of 7}
    \vskip 0.4cm
    {\color{white}\LARGE\bfseries Review, Quality Assurance, and Finalization}
    \vskip 0.4cm
    {\color{nsoLight}\large Building a multi-stage QA process that protects NSO Nepal's institutional credibility.}
    \vskip 0.8cm
    \textcolor{nsoGold}{\rule{3cm}{3pt}}
\end{frame}
}

% ==========================================
% WHY QA IS NON-NEGOTABLE
% ==========================================
\begin{frame}{Why QA Is Non-Negotiable}
    \begin{blockquote}
    A statistical report published with \textbf{one error} --- a wrong number, a mislabelled table, an inconsistency between text and chart --- damages the credibility of the institution, not just the author.
    \end{blockquote}
    \vskip 0.2cm
    For NSO Nepal --- \textbf{custodian of Nepal's official statistics} --- every report represents the office's reputation.
    \vskip 0.2cm
    \textbf{Types of errors to catch:}
    \begin{itemize}
        \item Numerical mismatch between text and tables.
        \item Wrong percentage/percentage point distinction.
        \item Inconsistent time period labelling.
        \item Mislabelled tables or figures.
        \item Executive summary inconsistent with findings.
    \end{itemize}
\end{frame}

% ==========================================
% QA CHECKLIST 1
% ==========================================
\begin{frame}{QA Checklist --- Factual \& Consistency}
    \begin{columns}[t]
        \begin{column}{0.48\textwidth}
            \textbf{Factual / Data Accuracy}
            \begin{itemize}
                \item[$\square$] All numbers in text match numbers in tables.
                \item[$\square$] Comparisons to previous periods are correct.
                \item[$\square$] Totals and sub-totals cross-checked.
                \item[$\square$] Percentage change calculations verified.
            \end{itemize}
        \end{column}
        \begin{column}{0.48\textwidth}
            \textbf{Consistency}
            \begin{itemize}
                \item[$\square$] Same number written the same way throughout.
                \item[$\square$] Table and figure numbers sequential \& correctly referenced.
                \item[$\square$] Time periods described consistently.
            \end{itemize}
            \vskip 0.2cm
            \textbf{Language and Clarity}
            \begin{itemize}
                \item[$\square$] Executive summary matches findings.
                \item[$\square$] All jargon defined on first use.
            \end{itemize}
        \end{column}
    \end{columns}
\end{frame}

% ==========================================
% QA CHECKLIST 2
% ==========================================
\begin{frame}{QA Checklist --- Formatting \& Press Release}
    \begin{columns}[t]
        \begin{column}{0.48\textwidth}
            \textbf{Formatting and Style}
            \begin{itemize}
                \item[$\square$] NSO Nepal house style applied (fonts, logos, page numbering).
                \item[$\square$] All tables/figures properly titled and sourced.
                \item[$\square$] Table of Contents is up to date.
                \item[$\square$] PDF named correctly \\ (e.g., \textit{NSO\_CPI\_Baisakh2082\_Final.pdf}).
            \end{itemize}
        \end{column}
        \begin{column}{0.48\textwidth}
            \textbf{Press Release --- Specific Checks}
            \begin{itemize}
                \item[$\square$] Headline accurately reflects the most important number.
                \item[$\square$] Dated and marked ``For Immediate Release''.
                \item[$\square$] Contact details included.
                \item[$\square$] Quote reviewed and approved by the senior official.
            \end{itemize}
        \end{column}
    \end{columns}
\end{frame}

% ==========================================
% THE REVIEW WORKFLOW
% ==========================================
\begin{frame}{The Review Workflow}
    \centering
    \begin{tikzpicture}[
        node distance=0.3cm and 0.5cm,
        wf/.style={rectangle, draw=nsoBlue, fill=nsoLight, text=nsoBlue, font=\scriptsize\bfseries, minimum width=3.2cm, minimum height=0.6cm, align=center, rounded corners=3pt}
    ]
        \node[wf] (a) {Author};
        \node[wf, right=of a] (p) {Subject-matter Peer Review \\ \tiny (checks facts, analysis)};
        \node[wf, right=of p] (t) {Technical Review \\ \tiny (checks methods, tables)};
        \node[wf, below=of t] (m) {Management Sign-off \\ \tiny (Chief Statistician)};
        \node[wf, left=of m] (pb) {Publication};

        \draw[->, >=stealth, nsoGold, line width=1pt] (a) -- (p);
        \draw[->, >=stealth, nsoGold, line width=1pt] (p) -- (t);
        \draw[->, >=stealth, nsoGold, line width=1pt] (t) -- (m);
        \draw[->, >=stealth, nsoGold, line width=1pt] (m) -- (pb);
    \end{tikzpicture}
    \vskip 0.4cm
    \begin{itemize}
        \item \textbf{For major publications} (census, survey reports), also consider:
        \begin{itemize}
            \item \textbf{External expert review} --- international consultants, academic peers.
            \item \textbf{Pre-publication stakeholder consultation} --- relevant ministries.
        \end{itemize}
    \end{itemize}
\end{frame}

% ==========================================
% FINALIZATION CHECKLIST
% ==========================================
\begin{frame}{Finalization Checklist --- Before Publishing}
    \rowcolors{2}{nsoLight}{white}
    \small
    \begin{tabular}{lc}
    \toprule
    \rowcolor{nsoBlue}\color{white}\bfseries Item & \color{white}\bfseries Status \\
    \midrule
    All data verified against source datasets & [ $\square$ ] \\
    Final text approved by Chief Statistician & [ $\square$ ] \\
    Press release prepared and ready to send simultaneously & [ $\square$ ] \\
    PDF locked and named correctly & [ $\square$ ] \\
    Report uploaded to NSO Nepal website with correct metadata & [ $\square$ ] \\
    Press release distributed to media list & [ $\square$ ] \\
    Social media posts scheduled (if applicable) & [ $\square$ ] \\
    \bottomrule
    \end{tabular}
    \vskip 0.2cm
    \begin{takeawaybox}
    \scriptsize \textbf{Key Takeaway:} QA is a process, not a single proofread. Build review stages into your production schedule as planned milestones --- not afterthoughts.
    \end{takeawaybox}
\end{frame}

% ==========================================
% STEP 7 HEADER
% ==========================================
{
\setbeamercolor{background canvas}{bg=nsoBlue}
\setbeamercolor{normal text}{fg=white}
\setbeamercolor{structure}{fg=white}
\begin{frame}[plain]
    \vspace{1.5cm}
    \stepbadge{Step 7 of 7}
    \vskip 0.4cm
    {\color{white}\LARGE\bfseries Synthesis: From Blank Page to Published Report}
    \vskip 0.4cm
    {\color{nsoLight}\large Applying all six steps to a real NSO Nepal publication scenario.}
    \vskip 0.8cm
    \textcolor{nsoGold}{\rule{3cm}{3pt}}
\end{frame}
}

% ==========================================
% WORKED EXAMPLE CPI
% ==========================================
\begin{frame}{Worked Example: CPI Baisakh 2082}
    \textbf{Situation:} End of Baisakh. Monthly CPI data processing is complete. You manage the publication from start to finish.
    \vskip 0.3cm
    \textbf{Stage 1 --- Identify what to produce (Step 2)}
    \vskip 0.2cm
    \rowcolors{2}{nsoLight}{white}
    \footnotesize
    \begin{tabular}{lll}
    \toprule
    \rowcolor{nsoBlue}\color{white}\bfseries Document & \color{white}\bfseries Type & \color{white}\bfseries Audience \\
    \midrule
    CPI Monthly Bulletin & Statistical Bulletin & Government, researchers, public \\
    Press Release & Press Release & Media, general public \\
    Internal Briefing Note & Administrative Report & Chief Statistician / NPC \\
    \bottomrule
    \end{tabular}
\end{frame}

% ==========================================
% STAGE 2 OUTLINE
% ==========================================
\begin{frame}{Stage 2 --- Build the Outline (Step 3)}
    \begin{columns}[t]
        \begin{column}{0.48\textwidth}
            \begin{tcolorbox}[colback=nsoLight, colframe=nsoBlue, title=CPI Bulletin Outline, coltitle=white, boxrule=1pt]
                \footnotesize
                Cover $\rightarrow$ Key Highlights Box $\rightarrow$ \\
                1. Introduction $\rightarrow$ 2. Overall CPI $\rightarrow$ \\
                3. Food Inflation $\rightarrow$ 4. Non-Food $\rightarrow$ \\
                5. Regional Breakdown $\rightarrow$ \\
                6. Methodology Note $\rightarrow$ Data Tables \\
                $\rightarrow$ Contact
            \end{tcolorbox}
        \end{column}
        \begin{column}{0.48\textwidth}
            \begin{tcolorbox}[colback=nsoLight, colframe=nsoBlue, title=Press Release Outline, coltitle=white, boxrule=1pt]
                \footnotesize
                Headline (overall CPI + change) $\rightarrow$ \\
                Opening paragraph (headline in plain English) $\rightarrow$ Food inflation $\rightarrow$ \\
                Regional highlight $\rightarrow$ Context $\rightarrow$ \\
                NSO Chief Quote $\rightarrow$ Boilerplate $\rightarrow$ Contact
            \end{tcolorbox}
        \end{column}
    \end{columns}
\end{frame}

% ==========================================
% STAGE 3 WRITE FOR AUDIENCE
% ==========================================
\begin{frame}{Stage 3 --- Write for Each Audience (Steps 4 \& 5)}
    \begin{columns}[t]
        \begin{column}{0.48\textwidth}
            \small
            \textbf{Statistical Bulletin}
            \begin{itemize}
                \item Full methodology note.
                \item Precise tables with source lines.
                \item Percentage point language throughout.
                \item Comparison to previous month \& year.
            \end{itemize}
            \vskip 0.2cm
            \textbf{Internal Briefing Note}
            \begin{itemize}
                \item Bullet points only.
                \item 1 page maximum.
                \item Summary + 2--3 recommendations.
                \item No methodology detail needed.
            \end{itemize}
        \end{column}
        \begin{column}{0.48\textwidth}
            \small
            \textbf{Press Release}
            \begin{itemize}
                \item One headline number up top.
                \item Plain language, no jargon.
                \item Round headline figures.
                \item One NSO Chief quote in plain language.
                \item Methodology in one sentence.
                \item ``For Immediate Release'' header.
            \end{itemize}
        \end{column}
    \end{columns}
\end{frame}

% ==========================================
% STAGE 4 REVIEW
% ==========================================
\begin{frame}{Stage 4 --- Review \& Publish (Step 6)}
    \centering
    \begin{tikzpicture}[
        node distance=0.4cm,
        stepbox/.style={rectangle, draw=nsoBlue, fill=nsoLight, text=nsoBlue, font=\scriptsize\bfseries, minimum width=9cm, minimum height=0.5cm, align=center, rounded corners=2pt}
    ]
        \node[stepbox] (s1) {Peer reviewer checks all figures match the source data};
        \node[stepbox, below=of s1] (s2) {Manager reviews press release quote accuracy};
        \node[stepbox, below=of s2] (s3) {Chief Statistician signs off};
        \node[stepbox, below=of s3, fill=nsoGold, draw=nsoGold, text=white] (s4) {10:00 AM — Publication Day: Bulletin + Press Release + Web update released simultaneously};

        \draw[->, >=stealth, nsoGold, line width=1pt] (s1) -- (s2);
        \draw[->, >=stealth, nsoGold, line width=1pt] (s2) -- (s3);
        \draw[->, >=stealth, nsoGold, line width=1pt] (s3) -- (s4);
    \end{tikzpicture}
    \vskip 0.3cm
    \begin{takeawaybox}
    \footnotesize \textbf{Note:} Simultaneous release is critical. If the press release goes out before the report is on the website, journalists cannot verify the numbers themselves.
    \end{takeawaybox}
\end{frame}

% ==========================================
% THE MENTAL MODEL
% ==========================================
\begin{frame}{The Mental Model --- Always}
    \begin{columns}[c]
        \begin{column}{0.48\textwidth}
            \centering
            \begin{tikzpicture}[scale=0.8, every node/.style={transform shape},
                node distance=0.15cm,
                flowbox/.style={
                    rectangle,
                    draw=none,
                    fill=nsoBlue,
                    text=white,
                    font=\scriptsize\bfseries,
                    text width=4.5cm,
                    minimum height=0.5cm,
                    align=center,
                    rounded corners=3pt
                },
                flowboxgold/.style={
                    flowbox,
                    fill=nsoGold
                },
                arrow/.style={
                    -{Stealth[scale=0.8]},
                    nsoGold,
                    line width=1.2pt
                }
            ]
                \node[flowbox] (p) {Purpose + Audience};
                \node[flowbox, below=of p] (r) {Report Type};
                \node[flowbox, below=of r] (o) {Outline};
                \node[flowbox, below=of o] (s) {Section Content};
                \node[flowbox, below=of s] (sp) {Statistical Precision};
                \node[flowbox, below=of sp] (q) {Quality Review};
                \node[flowboxgold, below=of q] (pb) {Publication};
                
                \draw[arrow] (p) -- (r);
                \draw[arrow] (r) -- (o);
                \draw[arrow] (o) -- (s);
                \draw[arrow] (s) -- (sp);
                \draw[arrow] (sp) -- (q);
                \draw[arrow] (q) -- (pb);
            \end{tikzpicture}
        \end{column}
        \begin{column}{0.48\textwidth}
            \small
            Every official report --- from a \textbf{2-page press release} to a \textbf{300-page census} --- follows the same logic.
            \vskip 0.4cm
            The \textbf{scale} changes. \\
            The \textbf{thinking} does not.
            \vskip 0.4cm
            \begin{blockquote}
            Master this workflow and you can confidently approach any official report at NSO Nepal.
            \end{blockquote}
        \end{column}
    \end{columns}
\end{frame}

% ==========================================
% SUMMARY
% ==========================================
\begin{frame}{Summary --- What You Now Know}
    \rowcolors{2}{nsoLight}{white}
    \scriptsize
    \begin{tabular}{cl}
    \toprule
    \rowcolor{nsoBlue}\color{white}\bfseries Step & \color{white}\bfseries Key Idea \\
    \midrule
    \textbf{1} & Reports differ by type --- identify the right one before writing \\
    \textbf{2} & Audience + purpose + occasion $\rightarrow$ determines report type; press releases accompany every major release \\
    \textbf{3} & Outline before you write --- it is a thinking tool, not just a contents page \\
    \textbf{4} & Every section has a specific job; press releases follow the inverted pyramid \\
    \textbf{5} & Anchor numbers in time, population, and units; distinguish \% from percentage points \\
    \textbf{6} & QA is multi-stage --- build it into your production schedule \\
    \textbf{7} & All steps form a repeatable workflow from blank page to published report \\
    \bottomrule
    \end{tabular}
\end{frame}

% ==========================================
% WHERE TO GO FROM HERE
% ==========================================
\begin{frame}{Where to Go From Here}
    \begin{columns}[t]
        \begin{column}{0.31\textwidth}
            \begin{tcolorbox}[colback=nsoLight, colframe=nsoBlue, title=1. International Standards, coltitle=white, boxrule=1pt, height=3.8cm]
                \scriptsize UN Statistics Division guidelines, IMF Data Standards (SDDS), and how they shape NSO Nepal norms.
            \end{tcolorbox}
        \end{column}
        \begin{column}{0.31\textwidth}
            \begin{tcolorbox}[colback=nsoLight, colframe=nsoBlue, title=2. Data Visualisation, coltitle=white, boxrule=1pt, height=3.8cm]
                \scriptsize Designing tables, charts, and infographics that meet accessibility and publication standards.
            \end{tcolorbox}
        \end{column}
        \begin{column}{0.31\textwidth}
            \begin{tcolorbox}[colback=nsoLight, colframe=nsoBlue, title=3. Plain Language, coltitle=white, boxrule=1pt, height=3.8cm]
                \scriptsize Making statistical findings understandable and accessible to non-technical audiences.
            \end{tcolorbox}
        \end{column}
    \end{columns}
\end{frame}

% ==========================================
% THANK YOU
% ==========================================
{
\setbeamercolor{background canvas}{bg=nsoBlue}
\setbeamercolor{normal text}{fg=white}
\setbeamercolor{structure}{fg=white}
\begin{frame}[plain]
    \vspace{1.5cm}
    \textcolor{nsoGold}{\rule{3cm}{4pt}}
    \vskip 0.4cm
    {\color{white}\huge\bfseries Thank You}
    \vskip 0.2cm
    {\color{nsoLight}\Large National Statistics Office Nepal} \\
    {\color{nsoMid}\large Government of Nepal}
    \vskip 0.4cm
    \textcolor{nsoGold}{\rule{3cm}{4pt}}
    \vskip 0.4cm
    \begin{columns}[t]
    \begin{column}{0.5\textwidth}
        {\color{white}\scriptsize\bfseries Planning and Structuring of Official Report Writing} \\[2pt]
        {\color{nsoLight}\tiny A Learning Guide for NSO Nepal Staff}
    \end{column}
    \end{columns}
    \vskip 0.8cm
    {\color{nsoMid}\tiny June 2026 · For Internal Distribution · NSO Nepal Staff}
\end{frame}
}

\end{document}